\documentclass[11pt, a4paper]{article}

% ==============================================================================
% PAQUETES Y CONFIGURACIÓN GENERAL
% ==============================================================================
\usepackage[spanish, es-tabla, es-nodecimaldot]{babel}
\usepackage[utf8]{inputenc}
\usepackage[T1]{fontenc}
\usepackage{amsmath, amsfonts, amssymb}
\usepackage{geometry}
\geometry{top=2.5cm, bottom=2.5cm, left=2.5cm, right=2.5cm}
\usepackage{xcolor}
\usepackage{enumitem}
\usepackage{booktabs}
\usepackage{longtable}
\usepackage{graphicx}
\usepackage{caption}
\usepackage{tcolorbox}
\tcbuselibrary{skins, breakable}

\usepackage{hyperref}
\hypersetup{
    colorlinks=true,
    linkcolor=blue!80!black,
    citecolor=blue!80!black,
    urlcolor=blue!80!black,
    pdfencoding=unicode
}

% ==============================================================================
% DEFINICIÓN DE ENTORNOS Y CAJAS EN LA PLANTILLA
% ==============================================================================
\newtcolorbox{definition}[1][]{%
    enhanced,
    colback=blue!5!white,
    colframe=blue!75!black,
    fonttitle=\bfseries,
    title={Definición: #1},
    breakable,
    arc=0mm
}

\newtcolorbox{important}[1][]{%
    enhanced,
    colback=red!5!white,
    colframe=red!75!black,
    fonttitle=\bfseries,
    title={Importante: #1},
    breakable,
    arc=0mm
}

\newtcolorbox{example}[1][]{%
    enhanced,
    colback=green!5!white,
    colframe=green!50!black,
    fonttitle=\bfseries,
    title={Ejemplo: #1},
    breakable,
    arc=0mm
}

\newtcolorbox{exercise}[1][]{%
    enhanced,
    colback=yellow!10!white,
    colframe=yellow!50!black,
    fonttitle=\bfseries,
    title={Ejercicio: #1},
    breakable,
    arc=0mm
}

% ==============================================================================
% DATOS DEL DOCUMENTO
% ==============================================================================
\title{\textbf{Derecho Registral y Derechos Reales}\\\large Apunte Académico de Estudio}
\author{\textbf{Material de Estudio Cátedra}}
\date{\today}

\begin{document}

\maketitle
\tableofcontents
\clearpage

% ==============================================================================
% DESARROLLO DEL CONTENIDO ACADÉMICO
% ==============================================================================

\section{Introducción a los Derechos Reales}

\begin{definition}[Derecho Real]
Un \textbf{derecho real} es el poder jurídico, de estructura legal, que se ejerce directamente sobre su objeto, en forma autónoma y que atribuye a su titular las facultades de persecución y preferencia, y las demás previstas en el Código Civil y Comercial.
\end{definition}

El régimen de los derechos reales se caracteriza por el principio de \textbf{orden público} y la regla del \textit{numerus clausus} (número cerrado), fijada por ley.

\subsection{Estructura y Elementos}

\begin{itemize}
    \item \textbf{Sujeto:} El titular del derecho real (persona humana o jurídica).
    \item \textbf{Objeto:} Cosas individuales y determinadas en el comercio, o bienes en los casos expressamente señalados por la ley.
    \item \textbf{Relación directa e inmediata:} El titular obtiene la utilidad de la cosa sin necesidad de la intervención de un tercero.
\end{itemize}

\begin{important}[Absolutidad y Oponibilidad]
Los derechos reales son \textbf{absolutos} y oponibles \textit{erga omnes} (frente a todos). De este carácter surge la necesidad insoslayable de la \textbf{publicidad registral} para que sean conocidos y respetados por terceros.
\end{important}

\section{Publicidad Registral y Sus Principios}

La \textbf{publicidad registral} es el conjunto de medios dirigidos a hacer saber a la comunidad las situaciones jurídicas relativas a la titularidad y cargas de los bienes inscriptos.

\subsection{Principios Registrales Fundamentales}

\begin{enumerate}
    \item \textbf{Principio de Inscripción:} Determina la exigencia de registrar un título para que tenga efectos legales plenamente oponibles.
    \item \textbf{Principio de Prioridad:} Establece que el acto inscripto en primer lugar prevalece sobre los posteriores. Se expresa mediante el aforismo \textit{prior tempore, potior iure}.
    \item \textbf{Principio de Tracto Sucesivo:} Exige que la serie de inscripciones mantenga una cadena ininterrumpida de titularidades. El disponente actual debe figurar previamente como titular inscripto.
    \item \textbf{Principio de Legalidad:} Requiere la calificación o control previo del documento formal por parte del registrador antes de proceder a la toma de razón.
    \item \textbf{Principio de Rogación:} Los asientos registrales se efectúan siempre a instancia de parte legítima o por orden judicial, nunca de oficio.
\end{enumerate}

\begin{example}[Aplicación del Principio de Prioridad]
Si el titular de un inmueble constituye dos hipotecas sobre el mismo bien en distintas fechas, la hipoteca inscripta en primer lugar ostentará prioridad cobratoria frente a la segunda en un eventual procedimiento judicial de ejecución.
\end{example}

\subsection{Efectos de la Inscripción Registral}

Los efectos de la inscripción registral de derechos reales sobre inmuebles se clasifican en:

\begin{itemize}
    \item \textbf{Efecto Declarativo:} Constituye la regla general del sistema inmobiliario argentino (Ley 17.801). El derecho real nace fuera del registro (mediante título y modo suficientes) y la inscripción se exige exclusivamente a los fines de la \textbf{oponibilidad a terceros interesados de buena fe}.
    \item \textbf{Efecto Constitutivo:} Caso excepcional en el cual el derecho real nace formalmente con la toma de razón en el registro respectivo (ejemplo: régimen automotor).
\end{itemize}

\begin{exercise}[Caso Práctico de Análisis]
Un comprador adquiere un inmueble mediante escritura pública y toma la posesión del bien, pero omite inscribir el título en el Registro de la Propiedad Inmueble. Un acreedor del vendedor traba un embargo sobre dicho inmueble con posterioridad a la compraventa.

\textbf{Consigna:} Determine la situación jurídica entre el comprador y el acreedor embargante aplicando los principios de publicidad declarativa y oponibilidad.
\end{exercise}

\clearpage

% ==============================================================================
% RESUMEN CORNELL — PREGUNTAS Y PISTAS DE REPASO
% ==============================================================================

\section{Resumen Cornell: Preguntas de Repaso}

\begin{description}[style=unboxed, leftmargin=0cm, labelindent=0pt]

    \item[\textbf{¿Qué es un derecho real y qué facultades atribuye a su titular?}]
    Es un poder jurídico ejercido de forma directa e inmediata sobre una cosa, que confiere a su titular las facultades de persecución y preferencia, además de ser oponible \textit{erga omnes}.

    \item[\textbf{¿Cuál es la diferencia sustancial entre efecto declarativo y constitutivo de la inscripción?}]
    En la inscripción declarativa el derecho real nace antes de su registro (por título y modo) y la inscripción le otorga oponibilidad frente a terceros. En la constitución registral, el derecho real nace únicamente a partir de su inscripción.

    \item[\textbf{¿En qué consiste el principio registral de tracto sucesivo?}]
    Es la exigencia legal de que cada asiento registral derive directamente del titular inscripto antecedente, garantizando un encadenamiento ininterrumpido de transmisiones.

    \item[\textbf{¿Qué función cumple el principio de legalidad o calificación registral?}]
    Obliga al registro a verificar formalmente los documentos presentados para constatar que cumplan los requisitos legales antes de autorizar su inscripción o anotación.

    \item[\textbf{¿Por qué los derechos reales requieren de la publicidad registral?}]
    Debido a su carácter absoluto y oponibilidad \textit{erga omnes}, es imprescindible que los terceros puedan conocer la existencia y situación de las cargas o titularidades sobre los bienes.

\end{description}

\end{document}